% part: intuitionistic-logic
% chapter: introduction
% section: natural-deduction

\documentclass[../../../include/open-logic-section]{subfiles}

\begin{document}

\olfileid{int}{int}{ntd}

\olsection{సహజ నిగమనం}

అంతఃప్రజ్ఞావాద తర్కానికి $\FalseCl$ నియమాలు లేని
సహజ నిగమనం ఒక ప్రామాణిక \tetoken{వ్యుత్పత్తి}{!!{derivation}}
వ్యవస్థ. నియమాలను ఇక్కడ మళ్లీ ఇచ్చి, BHK
అర్థనిర్దేశం ద్వారా వాటి ప్రేరణను సూచిస్తాం.
ప్రతి సందర్భంలోనూ పూర్వాపేక్షలకు నిర్మాణాలు
ఉంటే ముగింపుకూ నిర్మాణం ఉంటుందని నిర్ణయించడానికి
అనుమతించే నియమంగా దీనిని చూడవచ్చు.

సహజ నిగమన \tetoken{వ్యుత్పత్తులకు}{!!{derivation}s}
ఉపసంహరించని పరికల్పనలు ఉంటాయి. కాబట్టి
\tetoken{ఉపసంహరించని}{!!{undischarged}}
పరికల్పనలు~$\Gamma$ నుంచి $!A$ యొక్క
\tetoken{ఒక వ్యుత్పత్తిని}{!!a{derivation}},
ప్రతి $!B \in \Gamma$ నిర్మాణాలను~$!A$
నిర్మాణంగా మార్చే ప్రమేయకంగా చూడాలి.
\tetoken{ఉపసంహరించని}{!!{undischarged}}
పరికల్పనలు ఏవీ లేకుండా~$!A$కు
\tetoken{ఒక వ్యుత్పత్తి}{!!a{derivation}} ఉంటే,
BHK అర్థనిర్దేశం ప్రకారం~$!A$కు నిర్మాణం
ఉంటుంది. అయితే ఈ చర్చలో అవసరం లేనప్పుడు
$\Gamma$ను పేర్కొనము.

పరికల్పన~$!A$ ఒక్కటే, \tetoken{ఉపసంహరించని}{!!{undischarged}}
పరికల్పన~$!A$ నుంచి~$!A$కు
\tetoken{ఒక వ్యుత్పత్తి}{!!a{derivation}}.
ఇది BHK అర్థనిర్దేశానికి అనుగుణం: నిర్మాణాలపై
తాదాత్మ్య ప్రమేయకం~$!A$ యొక్క ఏ నిర్మాణాన్నైనా
తిరిగి~$!A$ నిర్మాణంగానే ఇస్తుంది.

\subsection{సంయోజనం}

\begin{defish}
\AxiomC{$!A$}
\AxiomC{$!B$}
\RightLabel{\Intro{\land}}
\BinaryInfC{$!A \land !B$}
\DisplayProof
\hfill
\begin{tabular}{l}
\AxiomC{$!A \land !B$}
\RightLabel{\Elim{\land}}
\UnaryInfC{$!A$}
\DisplayProof
\\[3ex]
\AxiomC{$!A \land !B$}
\RightLabel{\Elim{\land}}
\UnaryInfC{$!B$}
\DisplayProof
\end{tabular}
\end{defish}

వరుసగా $!A_1$, $!A_2$ నిర్మాణాలు $N_1$, $N_2$
మనకు ఉన్నాయనుకుందాం. అప్పుడు $!A_1 \land !A_2$
నిర్మాణం కూడా ఉంది; అదే జత $\tuple{N_1, N_2}$.

BHK అర్థనిర్దేశంలో $!A_1 \land !A_2$ నిర్మాణం
జత $\tuple{N_1, N_2}$. కాబట్టి అలాంటి జత
ఉందనుకుందాం. అప్పుడు ప్రతి సంయోజ్యానికీ
నిర్మాణం ఉంది: $N_1$ అనేది~$!A_1$ నిర్మాణం,
$N_2$ అనేది $!A_2$ నిర్మాణం.
\sourcecorrection{OLTEINTND-001}{మూలం ఈ వాక్యంలో సంయోగాన్ని $!A_1 \land !A_1$గా వ్రాసింది; అదే వాక్యంలోని $N_2$ నిర్మాణం, పూర్వ సంయోగం ప్రకారం రెండో సంయోజ్యాన్ని $!A_2$గా సరిచేశాం.}

\subsection{సోపాధికం}

\begin{defish}
  \AxiomC{$\Discharge{!A}{u}$}
  \DeduceC{$!B$}
  \DischargeRule{$\Intro{\lif}$}{u}
  \UnaryInfC{$!A \lif !B$}
  \DisplayProof
\hfill
  \AxiomC{$!A \lif !B$}
  \AxiomC{$!A$}
  \RightLabel{$\Elim{\lif}$}
  \BinaryInfC{$!B$}
  \DisplayProof
\end{defish}

\tetoken{ఉపసంహరించని}{!!{undischarged}}
పరికల్పన~$!A$ నుంచి~$!B$కు
\tetoken{ఒక వ్యుత్పత్తి}{!!a{derivation}} ఉంటే,
$!A$ నిర్మాణాలను~$!B$ నిర్మాణాలుగా మార్చే
ప్రమేయకం~$f$ ఉంటుంది. అదే ప్రమేయకం $!A \lif !B$
నిర్మాణం. కాబట్టి $\Intro\lif$ నియమపు పూర్వాపేక్షకు
$!A$ నిర్మాణంపై ఆధారపడిన నిర్మాణం ఉంటే,
ముగింపు $!A \lif !B$కూ నిర్మాణం ఉంటుంది.

మరోవైపు, $!A$ నిర్మాణం~$N$, $!A \lif !B$
నిర్మాణం~$f$ ఉన్నాయనుకుందాం. $!A \lif !B$
నిర్మాణం $!A$ నిర్మాణాలను~$!B$ నిర్మాణాలుగా
మార్చే ప్రమేయకం. అందువల్ల $f(N)$ ఒక~$!B$
నిర్మాణం; అంటే $\Elim\lif$ ముగింపుకు నిర్మాణం
ఉంది.

\subsection{విడియోజనం}

\begin{defish}
\begin{tabular}{l}
\AxiomC{$!A$}
\RightLabel{\Intro{\lor}}
\UnaryInfC{$!A \lor !B$}
\DisplayProof
\\[3ex]
\AxiomC{$!B$}
\RightLabel{\Intro{\lor}}
\UnaryInfC{$!A \lor !B$}
\DisplayProof
\end{tabular}
\hfill
\AxiomC{$!A \lor !B$}
\AxiomC{$\Discharge{!A}{n}$}
\DeduceC{$!C$}
\AxiomC{$\Discharge{!B}{n}$}
\DeduceC{$!C$}
\DischargeRule{\Elim{\lor}}{n}
\TrinaryInfC{$!C$}
\DisplayProof
\end{defish}

$!A_i$ నిర్మాణం~$N_i$ ఉంటే, దాన్ని $!A_1 \lor !A_2$
నిర్మాణం~$\tuple{i, N_i}$గా మార్చవచ్చు.
మరోవైపు, $!A_1 \lor !A_2$ నిర్మాణం ఉందనుకుందాం;
అంటే $N_i$ అనేది~$!A_i$ నిర్మాణమైన జత
$\tuple{i, N_i}$ ఉంది. అలాగే వరుసగా $!A_1$,
$!A_2$ నిర్మాణాలను~$!C$ నిర్మాణాలుగా మార్చే
ప్రమేయకాలు $f_1$, $f_2$ ఉన్నాయనుకుందాం.
అప్పుడు $f_i(N_i)$ అనేది~$\Elim\lor$
ముగింపు~$!C$కు నిర్మాణం.

\subsection{అసంభవం}

\begin{defish}
  \AxiomC{$\lfalse$}
  \RightLabel{\FalseInt}
  \UnaryInfC{$!A$}
  \DisplayProof
\end{defish}

\tetoken{ఉపసంహరించని}{!!{undischarged}}
పరికల్పనలు~$!B_1$, \dots, $!B_n$ నుంచి~$\lfalse$కు
\tetoken{ఒక వ్యుత్పత్తి}{!!a{derivation}} ఉంటే,
$!B_1$, \dots, $!B_n$ నిర్మాణాలను~$\lfalse$
నిర్మాణంగా మార్చే ప్రమేయకం~$f(M_1,
\dots, M_n)$ ఉంటుంది. $\lfalse$కు నిర్మాణం
లేదు కాబట్టి $!B_1$, \dots, $!B_n$ అన్నింటికీ
ఒకేసారి నిర్మాణాలు ఉండలేవు. అందువల్ల $f$కు
ఈ ధర్మం కూడా ఉంటుంది: వరుసగా $M_1$, \dots,
$M_n$లు $!B_1$, \dots, $!B_n$ నిర్మాణాలు
\emph{అయితే}, $f(M_1, \dots, M_n)$ ఒక~$!A$
నిర్మాణం \emph{అవుతుంది}.

\subsection{$\lnot$ కోసం నియమాలు}

$\lnot !A$ను $!A \lif \lfalse$గా నిర్వచించాం
కాబట్టి ఖచ్చితంగా చెప్పాలంటే~$\lnot$కు నియమాలు
అవసరం లేదు. అయినా ఇస్తే అవి ఇలా ఉంటాయి:

\begin{defish}
\AxiomC{$\Discharge{!A}{n}$}
\noLine
\DeduceC{$\lfalse$}
\DischargeRule{\Intro{\lnot}}{n}
\UnaryInfC{$\lnot !A$}
\DisplayProof
\hfill
\AxiomC{$\lnot !A$}
\AxiomC{$!A$}
\RightLabel{\Elim{\lnot}}
\BinaryInfC{$\lfalse$}
\DisplayProof
\end{defish}


\subsection{\usetoken{వ్యుత్పత్తుల}{derivation} ఉదాహరణలు}

\begin{enumerate}
\item $\Proves !A \lif (\lnot !A \lif \lfalse)$, అంటే $\Proves !A
  \lif ((!A \lif \lfalse) \lif \lfalse)$
  \begin{prooftree}
    \AxiomC{$\Discharge{!A}{2}$}
    \AxiomC{$\Discharge{!A \lif \lfalse}{1}$}
    \RightLabel{\Elim\lif}
    \BinaryInfC{$\lfalse$}
    \DischargeRule{\Intro\lif}{1}
    \UnaryInfC{$(!A \lif \lfalse)\lif \lfalse$}
    \DischargeRule{\Intro\lif}{2}
    \UnaryInfC{$!A \lif (!A \lif \lfalse) \lif \lfalse$}
  \end{prooftree}

\item $\Proves ((!A \land !B) \lif !C) \lif (!A \lif (!B \lif !C))$
  \begin{prooftree}
    \AxiomC{$\Discharge{(!A \land !B) \lif !C}{3}$}
    \AxiomC{$\Discharge{!A}{2}$}
    \AxiomC{$\Discharge{!B}{1}$}
    \RightLabel{\Intro\land}
    \BinaryInfC{$!A \land !B$}
    \RightLabel{\Elim\lif}
    \BinaryInfC{$!C$}
    \DischargeRule{\Intro\lif}{1}
    \UnaryInfC{$!B \lif !C$}
    \DischargeRule{\Intro\lif}{2}
    \UnaryInfC{$!A \lif (!B \lif !C)$}
    \DischargeRule{\Intro\lif}{3}
    \UnaryInfC{$((!A \land !B) \lif !C) \lif (!A \lif (!B \lif !C))$}
  \end{prooftree}

\item $\Proves \lnot(!A \land \lnot !A)$, అంటే $\Proves (!A \land (!A
  \lif \lfalse))\lif \lfalse$
  \begin{prooftree}
    \AxiomC{$\Discharge{!A \land (!A \lif \lfalse)}{1}$}
    \RightLabel{\Elim\land}
    \UnaryInfC{$!A \lif \lfalse$}
    \AxiomC{$\Discharge{!A \land (!A \lif \lfalse)}{1}$}
    \RightLabel{\Elim\land}
    \UnaryInfC{$!A$}
    \RightLabel{\Elim\lif}
    \BinaryInfC{$\lfalse$}
    \DischargeRule{\Intro\lif}{1}
    \UnaryInfC{$(!A \land (!A \lif \lfalse))\lif \lfalse$}
  \end{prooftree}

\item $\Proves \lnot\lnot(!A \lor \lnot !A)$, అంటే $\Proves ((!A \lor
  (!A \lif \lfalse))\lif \lfalse)\lif \lfalse$
  \begin{prooftree}\hskip-3em
    \AxiomC{$\Discharge{(!A \lor (!A \lif \lfalse))\lif \lfalse}{2}$}
    \AxiomC{$\Discharge{(!A \lor (!A \lif \lfalse))\lif \lfalse}{2}$}
    \AxiomC{$\Discharge{!A}{1}$}
    \RightLabel{\Intro\lor}
    \UnaryInfC{$!A \lor (!A \lif \lfalse)$}
    \RightLabel{\Elim\lif}
    \BinaryInfC{$\lfalse$}
    \DischargeRule{\Intro\lif}{1}
    \UnaryInfC{$!A \lif \lfalse$}
    \RightLabel{\Intro\lor}
    \UnaryInfC{$!A \lor (!A \lif \lfalse)$}
    \RightLabel{\Elim\lif}
    \insertBetweenHyps{\hskip -4em}
    \BinaryInfC{$\lfalse$}
    \DischargeRule{\Intro\lif}{2}
    \UnaryInfC{$((!A \lor (!A \lif \lfalse))\lif \lfalse)\lif \lfalse$}
  \end{prooftree}
\end{enumerate}

\begin{prop}
  అంతఃప్రజ్ఞావాద తర్కపు సహజ నిగమనంలో
  $\Gamma \Proves !A$ అయితే, సాంప్రదాయిక తర్కంలోనూ
  $\Gamma \Proves !A$. ముఖ్యంగా $!A$ ఒక
  అంతఃప్రజ్ఞావాద సిద్ధాంతమైతే, అది సాంప్రదాయిక
  సిద్ధాంతం కూడా.
\end{prop}

\begin{proof}
  ప్రతి సహజ నిగమన నియమం సాంప్రదాయిక సహజ
  నిగమనంలోనూ నియమమే. కాబట్టి అంతఃప్రజ్ఞావాద
  తర్కంలోని ప్రతి \tetoken{వ్యుత్పత్తి}{!!{derivation}}
  సాంప్రదాయిక తర్కంలోనూ
  \tetoken{ఒక వ్యుత్పత్తి}{!!a{derivation}}.
\end{proof}

\begin{prob}
  కింది \tetoken{సూత్రాలకు}{!!{formula}s}
  అంతఃప్రజ్ఞావాద తర్కంలో \tetoken{వ్యుత్పత్తులు}{!!{derivation}s} ఇవ్వండి:
  \begin{enumerate}
    \item $(\lnot !A \lor !B) \lif (!A \lif !B)$
    \item $\lnot\lnot\lnot !A \lif \lnot !A$
    \item $\lnot\lnot (!A \land !B) \liff (\lnot\lnot !A\land \lnot \lnot !B)$
    \item $\lnot (!A \lor !B) \liff (\lnot !A \land \lnot !B)$
    \item $(\lnot !A \lor \lnot !B) \lif \lnot (!A \land !B)$
    \item $\lnot\lnot ( !A \land !B) \lif  (\lnot\lnot !A \lor
    \lnot\lnot !B)$
  \end{enumerate}
\end{prob}

\end{document}
